% Proof-of-Control — arXiv preprint
% Build: tectonic main.tex  (XeTeX engine; fetches packages automatically)
% STATUS: working draft circulated for co-author review and consent prior to submission.
\documentclass[11pt]{article}

\usepackage[margin=1in]{geometry}
\usepackage{tikz}
\usepackage{graphicx}
\usepackage{booktabs}
\usepackage{array}
\usepackage{amsmath, amssymb}
\usepackage{listings}
\usepackage[outputdir=.]{minted}
\usemintedstyle{friendly}
\definecolor{codepanel}{HTML}{F7F5F1}
\usepackage{xcolor}
\usepackage{colortbl}
\usepackage[hidelinks]{hyperref}
\usepackage{enumitem}
\usepackage{microtype}
\usepackage{caption}
\usepackage{siunitx}
\usepackage{amsthm}
\newtheorem{definition}{Definition}
\newtheorem{theorem}{Theorem}
\newtheorem{proposition}{Proposition}
\usepackage[section]{placeins}

% Advanced AI Society brand palette
\definecolor{aailime}{HTML}{CFFF04}
\definecolor{aailimetint}{HTML}{F2FFB8}
\definecolor{aaiink}{HTML}{0A0A0A}
\definecolor{aaired}{HTML}{FFE3E4}
\definecolor{aaineutral}{HTML}{E7E3DD}
\definecolor{aaipaper}{HTML}{F0EDEA}
\definecolor{aaimuted}{HTML}{8F8A82}
\definecolor{aaidarklime}{HTML}{7A9900}

\lstset{basicstyle=\ttfamily\footnotesize, breaklines=true, frame=single,
        framerule=0.3pt, xleftmargin=1em, columns=fullflexible}
\captionsetup{font=small, labelfont=bf}

\newcommand{\iat}{\texttt{iat}}
\newcommand{\poc}{Proof-of-Control}

% The Society's bracket mark, drawn natively (lime on the dark band)
\newcommand{\aaimark}[1]{%
\begin{tikzpicture}[x=#1, y=#1, line cap=butt]
  \fill[aailime] (0,0) rectangle (0.16,1);
  \fill[aailime] (0.16,0) rectangle (0.42,0.24);
  \fill[aailime] (0.16,0.76) rectangle (0.42,1);
  \fill[aailime] (1.04,0) rectangle (1.20,1);
  \fill[aailime] (0.78,0) rectangle (1.04,0.24);
  \fill[aailime] (0.78,0.76) rectangle (1.04,1);
\end{tikzpicture}}

\begin{document}

% ---------------------------------------------------------------- cover page
\thispagestyle{empty}

\begin{tikzpicture}[remember picture, overlay]
  \fill[aaiink] (current page.north west) rectangle
    ([yshift=-6.35cm]current page.north east);
  \fill[aailime] ([yshift=-6.35cm]current page.north west) rectangle
    ([yshift=-6.53cm]current page.north east);
  \node[anchor=north west, inner sep=0pt]
    at ([xshift=1in, yshift=-0.95cm]current page.north west) {\aaimark{0.62cm}};
  \node[anchor=north west, inner sep=0pt, text width=14.2cm]
    at ([xshift=1in+1.15cm, yshift=-0.98cm]current page.north west)
    {\sffamily\bfseries\footnotesize\color{aailime}%
     ADVANCED AI SOCIETY\; \textcolor{aaimuted}{$\cdot$}\; PROOF-OF-CONTROL
     INITIATIVE\\[3pt]
     {\normalfont\sffamily\footnotesize\color{aaimuted}
      Working draft v0.1 $\cdot$ August 2026 $\cdot$ for co-author review
      prior to submission}};
  \node[anchor=north west, inner sep=0pt, text width=\textwidth]
    at ([xshift=1in, yshift=-2.55cm]current page.north west)
    {\sffamily\bfseries\fontsize{21.5}{26}\selectfont\color{aaipaper}%
     Proof-of-Control: An Open Standard for Runtime Verifiability and
     Cryptographic Oversight in Autonomous AI Execution};
  \node[anchor=north west, inner sep=0pt, text width=\textwidth]
    at ([xshift=1in, yshift=-5.55cm]current page.north west)
    {\sffamily\small\color{aaipaper}%
     \textcolor{aailime}{\rule[0.5ex]{1.6em}{2pt}}\hspace{0.6em}%
     Prove, rather than promise, what your agents did.};
\end{tikzpicture}

\vspace*{5.9cm}

% Compact author block
{\small\noindent
\textbf{Jim Schwoebel}\textsuperscript{1,2} \quad
\textbf{Ken Huang}\textsuperscript{2,3} \quad
\textbf{Tricia Wang}\textsuperscript{2} \quad
\textbf{Michael Casey}\textsuperscript{2}

\smallskip\noindent
{\footnotesize
\textsuperscript{1}\,Quome \quad
\textsuperscript{2}\,Advanced AI Society, \url{https://advancedaisociety.org} \quad
\textsuperscript{3}\,Cloud Security Alliance\footnote{Prepared within the
Proof-of-Control Initiative of the Advanced AI Society. Authorship is limited to
contributors to this manuscript; the Initiative's working groups, board, and
advisory board are credited in the Acknowledgments. Additional co-authors will be
added only upon explicit consent and substantive contribution.}}}

\smallskip\noindent
{\footnotesize Specification, reference implementation, and all artifacts:
\url{https://github.com/AAI-Society/ov-poc-standard}}

\vspace{0.9em}
\noindent\textcolor{aailime}{\rule{\textwidth}{1.6pt}}
\vspace{0.2em}

\begin{center}
{\sffamily\bfseries\small ABSTRACT}
\end{center}
\vspace{-0.5em}
\noindent{\small
Suppose you let a piece of software act for you: read your records, call your
tools, move your money. Afterward it tells you what it did. How would you
check? Today, in nearly every deployment, you cannot. The only account of what
happened is the system's own, and the party writing that account is the party
you would be checking. We call this the \emph{Verifiability Gap}, and this
paper is about closing it. We describe \poc{}, an open standard with six
domains of verification covering 127 testable requirements; four tiers that
grade evidence by \emph{who you must trust to believe it}, with a hard line
between the tiers where you still trust somebody and the tiers where you do
not; and an architecture that intercepts every consequential action, decides
it inside a protected environment, and emits a signed record chained to
everything before it. We prove that this arrangement does \emph{not} do what
you would assume it does---an attester that signs ``I evaluated this
snapshot'' says nothing about what the machine actually went on to do---and we
show what you have to add to fix it. Then we built the thing and measured it.
The software costs \SI{201}{\micro\second} per action, about one percent of
our budget. Checking the path an agent took costs the same whether the path is
ten steps or five thousand. Nine attacks work against the naive design and fail
against ours---and three of them were defects in our own implementation,
found by review rather than by design. We also
measure what we did not expect: a monitor that watches the whole path refuses
42\% of perfectly legitimate work unless you give it an explicit place to
forget, which is exactly what people who study information flow have been
saying for fifty years. Because an auditor should be able to check one action
without downloading a million, we adopt Certificate Transparency's tree and
show it turns a \SI{123}{\mega\byte} download into a 544-byte proof. Because
evidence has to be believable long after it is written, we measure surviving
quantum computers: the hash chain already does, and FIPS~204 ML-DSA signatures
cost a factor of 38 in size---a storage bill, not a cryptography problem, and
one that halves again if you encode the evidence as CBOR instead of JSON. And
because none of this is worth anything if two implementations disagree about
which bytes they signed, we publish a machine-readable schema, a canonical
form, and signed test vectors.}

\clearpage

\section{The Problem}
\label{sec:intro}

A few years ago software did what you told it, step by step, and if you wanted
to know what it had done you read the log. Now we have agents. You give one an
objective in a sentence, and it decides for itself what to do: which database
to query, which API to call, what to try next when the first thing fails,
which sub-agent to hand a piece of the job to. It does this at machine speed,
and it does not do the same thing twice~\cite{paloalto2026,airia2026}.

That is genuinely useful. It also breaks something we have quietly relied on
for a long time, which is the ability to find out afterward what happened.

Here is the situation exactly. An agent reads a customer record, calls three
tools, and issues a refund. It reports success. Did it open only the records
it was allowed to open, or did it look at others while it was in there? Did it
stay inside the authority you gave it, or did something it read along the way
talk it into doing more? You cannot tell. There is a log, but the machine that
performed the actions is the machine that wrote the log, and if that machine
is compromised---or if the company running it would simply rather you did not
know---the log says whatever they want it to say. An agent reporting that it
behaved is not evidence that it behaved.

We call this the \emph{Verifiability Gap}: the absence of independent evidence
of what an AI system did~\cite{pocspec}. It is not a policy gap. Nobody is
confused about what the agent \emph{should} do. It is an evidence gap---and
those are different problems with different fixes.

The reason it matters more every year is arithmetic. As machines get cheaper
at doing things, they do more things, and the number of actions somebody ought
to check grows with them. But the checking is still done by people, and there
are the same number of people as before. Catalini, Hui and Wu work this out
carefully~\cite{catalini2026}: the cost of doing falls toward zero while the
cost of verifying is pinned to human attention. So verification becomes the
scarce thing. If you want to check machine work at machine scale, the checking
has to be done by machines too, and the evidence has to be the kind a machine
can check.

\paragraph{What we did.} We wrote a standard, we analyzed it, we found holes
in our own design---two by proving theorems and a third by building the code---
we fixed them, and then we measured what the whole thing costs. Specifically:

\begin{enumerate}[itemsep=1pt,topsep=2pt]
\item \textbf{The standard} (Section~\ref{sec:standard}): six domains of
verification with 127 testable requirements, four evidence properties, four
tiers that grade evidence by how much trust it demands, and three stages of
how thoroughly a claim was checked.
\item \textbf{The architecture} (Section~\ref{sec:arch}): interception of
every consequential action, decided inside a hardware-protected environment,
under the IETF attestation architecture~\cite{rfc9334}, emitting signed
tokens~\cite{rfc9711} chained together in the manner of tamper-evident
logs~\cite{merkle1987,crosby2009,rfc6962}.
\item \textbf{The analysis} (Sections~\ref{sec:formal}--\ref{sec:security}):
a precise statement of what an attacker can do, four properties we want, and
proofs. Two of the four do \emph{not} follow from the obvious design. We show
exactly what has to be added.
\item \textbf{The comparison} (Section~\ref{sec:gap}): how much of this eight
existing frameworks already ask for, coded requirement by requirement, with an
honest account of what that measurement can and cannot tell you---including a
\emph{cross-model audit} of our own coding
(Section~\ref{sec:crossmodel}), because the party that wrote the requirements
should not be the only party deciding whether anyone else already covers them.
The procedure and its acceptance rule are reported in full, and we recommend
it for any coverage mapping where the author and the coder are the same.
\item \textbf{The interoperable claim set} (Section~\ref{sec:interop}): a
machine-readable schema in two renderings, a canonical serialization that fixes
exactly which bytes a signature covers, and signed test vectors including ten
negative cases. Without these, two conforming implementations produce different
digests for the same action and the central binding property fails silently.
\item \textbf{The implementation} (Section~\ref{sec:eval}): working code, an
attack harness that runs eleven attacks with and without each defense,
measurements of what it all costs, what it takes to make the evidence outlive
the signature scheme that protects it, and---in
Section~\ref{sec:merkle}---logarithmic proofs that let an auditor check one
action out of millions without downloading the rest.
\end{enumerate}

\paragraph{One finding we did not go looking for.} Building those proofs broke
our own verifier. An operator holding the signing key can rewrite a record in
the middle of the log, recompute every link after it, and re-sign the lot; the
result replays perfectly, because nothing about it is internally inconsistent.
Our verifier compared the published step \emph{count} and never the published
\emph{root}, so it accepted the forgery. That is attack A9, and the fix is now
requirement C7.3.5. It is worth noting how it was found: not by proving
anything, but by implementing a performance optimization carefully enough to
have to say what a published root was supposed to guarantee.

\paragraph{One thing we could not do.} We did not have a machine with a real
hardware enclave, so in our implementation the protected environment is a
software stand-in: the signing key is never handed out, the ``measurement'' is
a digest of the policy code, and every step of the protocol runs for real. But
hardware isolation is not there. That means our timings leave out the cost of
crossing into a real enclave, and the true number will be larger. We say so
rather than guessing at a number we did not measure.

\section{Why the Obvious Answers Don't Work}
\label{sec:background}

Before describing what we did, it is worth being clear about why the things
people already do are not enough. Not because they are bad---they are
useful---but because each one answers a different question.

\subsection{The path is the problem}
\label{sec:path}

Old software followed a decision tree you could read in advance. An agent does
not. Given one objective it invents its own sequence of steps, reads whatever
comes back, and changes its mind. So you cannot check compliance before it
runs; there is nothing fixed to check.

Worse, you cannot check it one step at a time either, and this is the part
people miss. Consider two actions. First, the agent reads a confidential
customer record---which it is allowed to do; that is its job. Second, it makes
an outbound network call---which it is also allowed to do; it has to talk to
partners. Look at either action alone and you see nothing wrong. Look at them
in sequence and you are watching data walk out the door.

This is not a new observation. It is the oldest observation in the field of
information flow~\cite{denning1976,goguen1982}. What is new is the setting:
benchmarks for tool-using agents show that defenses which work fine on
isolated prompts fall apart over multi-step
sequences~\cite{agentdojo,agentharm}, that content retrieved from the world
hijacks agents reliably~\cite{greshake2023,injecagent}, and that harm usually
arrives as a series of individually reasonable calls~\cite{toolemu}. We
return to the information-flow connection in Section~\ref{sec:ifc}, because it
turns out to predict our most important measurement.

\subsection{Three things people do, and what each one misses}

\paragraph{Telling the model to behave.} You write a system prompt: do not do
X, always do Y. The trouble is that the instruction lives in the same stream
of text as everything the agent reads---user input, retrieved documents, tool
output. There is no wall between ``what I was told by my owner'' and ``what I
read on a web page.'' So a web page can tell it something else. And even
setting attacks aside, a prompt is a suggestion with statistical force, not a
barrier. If the model steps over the line, nothing stops it. There is no line;
there is only a preference.

\paragraph{Permissions.} Give the agent credentials for exactly the endpoints
it needs. This is real and you should do it. But permissions are blind to
context and to history~\cite{sp800207}. An agent that legitimately needs to
read the customer database, and legitimately needs to send email, has
everything it needs to exfiltrate the customer database by email, and every
individual call is authorized. That is the path problem again, and a
permission system cannot see it because it looks at one call at a time.

\paragraph{Guardrails inside the application.} Put the checks in the agent's
own code. Now ask: what happens when the agent is running generated code, or
has a shell tool? The checks are running with the same authority as the thing
they are checking. A process can skip its own checks. And the log it writes is
its own log~\cite{netbankaudit2026}. This is the deep problem with all three:
\emph{the thing being checked is doing the checking}.

\subsection{Why looking afterward is too late}

The usual backstop is the audit: collect the logs, look later. Two things have
gone wrong with that.

First, the time available has collapsed. An instruction hidden in a web page
can redirect an agent instantly, and the money moves before anyone reads
anything~\cite{airia2026}. Afterward is measured in weeks; the event was
measured in milliseconds.

Second---and this is the fatal one---ordinary logs are not evidence. They are
a story the system tells about itself. Anyone who controls the machine can
edit them, add to them, or quietly leave things out. So the audit rests on
trusting the operator, which is precisely the thing you were trying not to
have to do.

What you want instead is evidence made \emph{by the mechanism, as the action
happens}, anchored to something the operator cannot
reach~\cite{redhat2026,eqty2026}.

\section{What Others Have Built}
\label{sec:related}

\paragraph{Governing agents at runtime.} The Agent Control Specification
defines where to put the hooks and how to carry policy between
systems~\cite{acs2026,acsms2026}, and CSA's AARM defines what to do at the
boundary: approve, modify, defer, deny~\cite{csaaarm}. That is the enforcement
half, and it is necessary. But the record it leaves is still the operator's
record. Enforcement decides what an agent may do; we are after evidence of
what it did.

\paragraph{Attestation.} The IETF's remote attestation
architecture~\cite{rfc9334,ratsarch} gives us the vocabulary and the roles,
entity attestation tokens give us the format~\cite{rfc9711,tcg2024}, there is
work on using several verifiers at once~\cite{multiverifier}, and hardware
enclaves plus verifiable-compute tooling give us somewhere to stand that the
host cannot reach into~\cite{redhat2026,eqty2026}. These prove things about
\emph{platforms and code}. We are extending them to say something about
\emph{what an agent did}.

\paragraph{Tamper-evident logs.} Chaining records so that changing one breaks
the rest goes back to Merkle~\cite{merkle1987} and was worked out properly for
logging by Crosby and Wallach~\cite{crosby2009}, then deployed at scale in
Certificate Transparency~\cite{rfc6962,rfc9162}. That literature also tells us
what such logs \emph{cannot} do on their own: an operator can show different
histories to different people, or simply stop showing you the recent
part. Which is why Certificate Transparency needs gossip and consistency
proofs~\cite{ctgossip,chase2016}, and why systems like CHAINIAC bring in
witnesses~\cite{chainiac}. We use those results in
Section~\ref{sec:security}; as far as we can tell, nobody has applied them to
agent evidence.

\paragraph{Reference monitors.} The idea that every security-relevant
operation must go through a mechanism that is tamper-proof, always invoked,
and small enough to verify is Anderson's~\cite{anderson1972}, and
``always invoked'' is Saltzer and Schroeder's complete
mediation~\cite{saltzer1975}. Section~\ref{sec:mediation} is really just this
old requirement, applied carefully, and finding that the obvious agent design
fails it.

\paragraph{What enclaves cannot promise.} Enclaves have been attacked, by
architectural analysis~\cite{costan2016} and by transient-execution
tricks~\cite{foreshadow}, and their trust ultimately roots in a chip vendor.
We therefore write those down as assumptions instead of pretending they are
free.

\paragraph{Standards and regulation.} NIST's risk framework~\cite{nistairmf}
and its generative profile~\cite{nist6001}, ISO/IEC 42001~\cite{iso42001} and
23894~\cite{iso23894}, SOC~2~\cite{soc2tsc}, the EU AI Act~\cite{euaiact},
OWASP AISVS~\cite{owaspaisvs}, MITRE ATLAS~\cite{mitreatlas}, and zero-trust
architecture~\cite{sp800207} all govern neighboring ground.
Section~\ref{sec:gap} measures exactly how much of what we ask for they
already ask for, using the coding method HAARF developed for healthcare
agents~\cite{haarf2026}.

\paragraph{Identity.} Decentralized identifiers~\cite{w3cdid}, verifiable
credentials~\cite{w3cvc2}, and supply-chain attestation
formats~\cite{slsa,c2pa} give us the pieces for saying which agent acted on
whose behalf, and where the model came from.

\section{The Standard}
\label{sec:standard}

Here is the whole idea in one sentence: \emph{for every control you claim, the
mechanism that enforces it must produce evidence, while it runs, that it held;
you must say how independently that evidence can be checked; and you must say
what you are still asking people to take on faith.} A system has \poc{} only
when its evidence reaches the top two tiers~\cite{pocspec}. We use MUST and
SHOULD in the RFC~2119 sense~\cite{rfc2119}.

\subsection{The six things worth having evidence about}
The requirements---127 of them, each phrased as something you can go and
check---are grouped into six domains:

\begin{enumerate}[itemsep=1pt,topsep=2pt]
\item \textbf{Provenance}: which model actually ran, and where it and its
inputs came from, tied together with signed manifests and hash
links~\cite{slsa,c2pa}.
\item \textbf{Privacy}: what data was read and written---proved without
handing over the data you are protecting, which is the awkward part.
\item \textbf{Portability}: what crossed which boundary, between companies,
jurisdictions, or clouds, with the evidence chain surviving the trip.
\item \textbf{Authorization}: what authority was granted, whether each action
stayed inside it, and whether delegation was valid---judged over the
\emph{path}, not one action at a time.
\item \textbf{Identity}: which agent, on whose behalf~\cite{w3cdid,w3cvc2}.
\item \textbf{Security}: whether the environment was intact, whether the
controls held, which tools were called, and who holds the keys.
\end{enumerate}

Four more chapters cut across all six: how evidence is generated and what
properties it must have, how it is graded, where in the stack it applies
(we use the seven MAESTRO layers~\cite{csamaestro}), and how a claim is
checked and what residual trust must be disclosed.

\subsection{Four tiers, and the line that matters}
\label{sec:tiers}

The grading question is not ``did you use cryptography?'' Cryptography is
everywhere and proves nothing by itself. The question is: \emph{who do I have
to trust to believe this?} (Figure~\ref{fig:ladder}.)

At Tier~1 you trust the operator: they say so. At Tier~2 you trust somebody
else---an auditor, a certificate authority, a chip vendor. Note that a signed
log is Tier~2: the signature proves the log was not edited after it was
written, and proves exactly nothing about whether it was true when it was
written. At Tier~3, there is nobody left to trust; anyone can check the
evidence with published tools. At Tier~4, checking is not optional---the
system cannot act unless the evidence exists.

The line falls between Tier~2 and Tier~3. Below it you have authenticated
paperwork. Above it you have evidence. We insist on a hard line rather than a
smooth scale for a practical reason: a buyer can ask a yes-or-no question.
``Does your system have \poc{}?'' is a question a procurement officer can put
in a contract. ``How far along the trust-minimization spectrum are you?'' is
not. Every certification that ever built a market did this---FIPS validated or
not, PCI compliant or not.

\begin{figure}[t]
\centering
\includegraphics[width=0.86\linewidth]{figures/tier-ladder.pdf}
\caption{The four tiers, graded by who you must trust. The hard line sits
between Tier~2 and Tier~3: below it, authenticated paperwork; above it,
evidence a stranger can check. Requirement levels 1--4 line up one-for-one
with the tiers.}
\label{fig:ladder}
\end{figure}

Separately from the evidence, we grade how carefully the \emph{claim} was
checked: the operator says so (Self-Declared), an accredited assessor checked
(Third-Party Assessed), or it is checked continuously as the system runs
(Continuously Monitored). And because you cannot govern what you have not
found, the declared inventory of agents must be reconciled against what
automatic discovery actually finds.

\section{How It Works}
\label{sec:arch}

The architecture is easier to describe than it sounds
(Figure~\ref{fig:arch}). Every time the agent is about to do something that
matters, we stop it. We write down exactly what it is proposing to do, in a
form that hashes to a fixed value. We hand that description to a protected
environment the host cannot reach into. That environment checks it against the
policy, decides, adds a link to a chain of everything decided so far, signs
the result, and hands back a token. Only then does the action go out. If the
evidence cannot be written, the action does not happen.

\begin{figure}[t]
\centering
\includegraphics[width=0.92\linewidth]{figures/reference-architecture.pdf}
\caption{The pipeline, and how it maps onto the IETF attestation roles. The
agent host is the Attester; the enclave is the Attesting Environment; the
token is Evidence; a Verifier checks it against published reference values
before a Relying Party will act.}
\label{fig:arch}
\end{figure}

\subsection{Where we stop the agent}
There are eight places, taken from the Agent Control
Specification~\cite{acs2026,acsms2026}: when a task starts; when context is
assembled (retrieval, memory); when a plan is formed; before a tool call;
after a tool result comes back; when memory is written; when a sub-task is
handed to another agent; and when the task finishes. At each one we build the
snapshot and get back one of four answers: \textsc{allow}, \textsc{deny},
\textsc{modify} (sanitize it and proceed), or \textsc{escalate} (ask a
human).

\begin{figure}[t]
\centering
\includegraphics[width=0.95\linewidth]{figures/evidence-flow.pdf}
\caption{The same boundary, seen from the evidence side. Nothing gets out
without a record, including the things that were refused.}
\label{fig:flow}
\end{figure}

\subsection{The roles, in the standard vocabulary}
Mapping onto the attestation architecture~\cite{rfc9334}: the machine running
the agent and the hooks is the \emph{Attester}; the layer holding the runtime,
the model context, and the proposed action is the \emph{Target Environment};
the enclave that measures the policy code, checks the bundle signature,
evaluates, and signs is the \emph{Attesting Environment}; the signed token is
\emph{Evidence}~\cite{rfc9711,tcg2024}; whoever compares that against
published reference values is a \emph{Verifier}; and whoever refuses to act
without a good answer is a \emph{Relying Party}. Chains of sub-agents are
handled with several verifiers~\cite{multiverifier}.

\subsection{The mistake we almost made}
\label{sec:mediation}

Now here is the interesting part, and it is the part we got wrong at first.

The enclave signs a statement. It is worth reading that statement carefully,
because it is narrower than it looks. What it says is: \emph{given this
description, the authorized policy said yes.} That is all. It does not say the
description was true. It does not say the machine then went and did that
thing.

So suppose the host is compromised. The host is what builds the description,
and the host is also what holds the credentials and the network connection. It
can hand the enclave a description of something harmless, collect a perfectly
genuine \textsc{allow} token, and then go do something else entirely. Every
signature checks out. The evidence is impeccable. And it is a lie.

We call this \emph{snapshot substitution}. It is not a clever attack; it is
what happens if you forget Anderson's rule that the monitor has to be
\emph{always invoked}, not merely present~\cite{anderson1972,saltzer1975}. You
have built a tamper-evident diary when you thought you were building a
reference monitor.

The fix is that the channel through which the effect actually happens has to
be inside the same trust boundary as the decision. Concretely, one of these
must be true, and the claim has to say which:

\begin{enumerate}[itemsep=1pt,topsep=2pt]
\item \textbf{The enclave makes the call itself.} It holds the credentials and
the connection; the host has no path to the endpoint.
\item \textbf{The enclave issues a ticket.} A single-use capability
cryptographically tied to the description and the target, which the far end
checks before doing anything. A request without a matching ticket is refused
where it lands.
\item \textbf{The path is attested.} Traffic can only leave through an
enforcement point that is itself attested and only admits requests carrying
matching evidence.
\end{enumerate}

Option~2 is the only one that survives a fully compromised host without
putting the network stack inside the enclave; options~1 and~3 shrink the
problem rather than removing it. If you do none of these, you may still claim
the lower tiers---but you must not claim your evidence describes what your
system did, because it does not.

\subsection{Chaining the records}
Each step makes a link:
\begin{equation}
L_t = H\!\left(L_{t-1} \,\|\, H(\sigma_t^{\mathrm{snap}}) \,\|\, d_t\right),
\qquad L_0 = H(\mathrm{seed}_{\mathrm{init}}),
\end{equation}
where $\sigma_t^{\mathrm{snap}}$ is the description, $d_t$ the decision, and
$H$ a hash nobody knows how to find collisions in. The enclave signs each
link. Numbering the steps means a missing one leaves a hole. Periodically the
running total is published somewhere public, off the critical path.

That much is standard~\cite{merkle1987,crosby2009}. It is also not enough, and
Section~\ref{sec:security} shows exactly how it fails---an operator can stop
publishing, or publish two different histories to two different people. Fixing
those needs bounded publishing intervals and people comparing notes, which is
what Certificate Transparency figured out~\cite{rfc6962,ctgossip,chase2016}.

\subsection{What a record actually contains}
\label{sec:schema}

Listing~\ref{lst:eat} is a real token, from the tool-call hook, for a payment
API request. It is an entity attestation token~\cite{rfc9711} with our claims
inside it. Table~\ref{tab:claims} says what every field means.

Three choices in there are worth pointing out. First, \emph{the data never
leaves}: the token carries a hash of the description, not the description, so
anyone can later confirm what was evaluated by hashing it, and the evidence
store never becomes a second copy of your customer database. Second,
\emph{history is load-bearing}: the step number and the running total tie this
record to every record before it, so quietly dropping an inconvenient step
breaks arithmetic rather than merely being impolite. Third, \emph{the rules are
pinned}: the bundle hash names the exact policy that produced this answer, so
``which rules were in force?'' is something you check, not something you
reconstruct from a change log and hope.

\paragraph{How you check one.} It is mechanical, and worth walking through
because there is nothing clever in it. (1)~Check the signature against the
platform's attestation roots. (2)~Compare the enclave measurement to the
published golden value---this is what tells you the code that decided was the
code you think decided. (3)~Check the policy bundle hash against the registry.
(4)~Check the nonce is fresh, so this is not an old token being waved at you
again. (5)~Check the step numbers run consecutively and the running totals
agree with the other tokens you have seen. (6)~Only now believe the verdict.
Every one of those steps uses published material. You never ask the operator
for anything. That is the entire point.

\begin{listing}[p]
\begin{minted}[bgcolor=codepanel, fontsize=\scriptsize, frame=lines,
framerule=0.4pt, breaklines]{json}
{
  "iss": "https://verifier.trusted-attestation.org",
  "iat": 1773532800,
  "nonce": "0x8f93e2a110c49b",
  "eat_profile": "https://standards.org/poc/v1",
  "poc_claims": {
    "agent_id": "did:web:enterprise.com:agents:finance-agent-04",
    "initiating_user": "user:auth0|84f12a9e",
    "agbom_digest": "sha256:7f8a...5f6a",
    "interception_point": "PRE_CALL_TOOL_INVOCATION",
    "step_index": 3,
    "merkle_root": "0x2c3d...2c3d",
    "policy_bundle_hash": "sha256:e3b0...b855",
    "target_resource": "api.bank.com/v1/payments",
    "canonical_snapshot_hash": "sha256:a4b3...a1b2",
    "verdict": "ALLOW"
  },
  "submods": {
    "hardware_attestation": {
      "platform": "INTEL_TDX",
      "mr_config": "0x1a2b3c4d5e6f...",
      "attestation_signature": "0x30440220..."
    }
  }
}
\end{minted}
\caption{A \poc{} Evidence token (EAT, JWT rendering) emitted at the
pre-call tool-invocation hook. Claim-by-claim definitions follow in
Table~\ref{tab:claims} below.}
\label{lst:eat}
\vspace{1.2em}
\centering\scriptsize
\begin{tabular}{@{}llp{8.1cm}@{}}
\toprule
\textbf{Claim} & \textbf{Type} & \textbf{Meaning} \\
\midrule
\texttt{iss} / \texttt{iat} & EAT & Token issuer and issuance timestamp. \\
\texttt{nonce} & EAT & Relying-party challenge; proves freshness and
prevents replay of an old token for a new action. \\
\texttt{eat\_profile} & EAT & The \poc{} EAT profile (and version) this
token conforms to. \\
\texttt{agent\_id} & PoC & Decentralized identifier of the agent instance
(C5); resolvable to its public keys. \\
\texttt{initiating\_user} & PoC & The principal on whose authority the
action runs---the human-to-agent binding (C5.1). \\
\texttt{agbom\_digest} & PoC & Digest of the Agent Bill of Materials: the
build, tools, and model manifest the agent runs from (C1). \\
\texttt{interception\_point} & PoC & Which lifecycle hook produced this
token (one of the eight in \S\ref{sec:arch}). \\
\texttt{step\_index} & PoC & Monotonic position in the execution path;
gaps expose omitted steps. \\
\texttt{chain\_head} & PoC & Running hash link over all steps so far;
what a verifier replaying the whole history recomputes. \\
\texttt{merkle\_root} & PoC & RFC 6962 tree head over the same leaves; what
an \emph{inclusion proof} reconstructs (\S\ref{sec:merkle}). \\
\texttt{tree\_size} & PoC & Number of leaves the root covers. An inclusion
proof means nothing without it. \\
\texttt{policy\_bundle\_hash} & PoC & Digest of the signed Rego bundle
that was evaluated---pins the exact policy in force. \\
\texttt{target\_resource} & PoC & The resource the proposed action
addressed (tool endpoint, table, device). \\
\texttt{canonical\_snapshot\_hash} & PoC & Commitment to the full
canonical state snapshot that was evaluated; discloses nothing by itself. \\
\texttt{verdict} & PoC & The binding decision: \textsc{allow},
\textsc{deny}, \textsc{modify}, or \textsc{escalate}. \\
\texttt{platform} / \texttt{mr\_config} & HW & TEE platform and enclave
measurement; compared against published reference values to prove the
authorized policy engine ran unmodified. \\
\texttt{attestation\_signature} & HW & Hardware-rooted signature over the
token by the enclave-bound key. \\
\bottomrule
\end{tabular}
\captionof{table}{The Evidence claim set. ``EAT'' rows are standard Entity
Attestation Token claims~\cite{rfc9711}; ``PoC'' rows are the \poc{}
profile; ``HW'' rows form the hardware-attestation submodule. Every field
here has a machine-readable definition, a canonical serialization, and signed
test vectors---see \S\ref{sec:interop}, because a table like this one is not
enough to make two implementations agree.}
\label{tab:claims}
\end{listing}

\subsection{Making two implementations agree}
\label{sec:interop}

Table~\ref{tab:claims} is prose in a grid. It is enough for a reader and not
enough for an implementer, and the gap between those two is where a security
property quietly dies.

Here is how. Theorem~\ref{thm:fidelity} works because the relying party
recomputes a digest over the request it has been asked to perform and compares
it to the digest the enclave committed to. Now suppose your implementation
serializes an action with the keys in the order they were inserted, and mine
sorts them. Same action, different bytes, different digest. Every signature
still verifies. Nothing throws an exception. The comparison just fails, for no
reason anyone can see from the outside---and the natural thing for a tired
engineer to do at that point is loosen the comparison until the two systems
interoperate. Which removes the very thing the comparison was there for.

So the claim set now comes with three artifacts rather than a table. There is a
CDDL definition, which is the form an IETF attestation profile is actually
written in, plus a JSON Schema for deployments that use JWTs. There is a
document that fixes the canonical form: key ordering, string escaping, digest
case, what ``absent'' means, and exactly which bytes a signature covers. And
there are signed test vectors, which are the part that does the real work.

We made two decisions worth repeating. Digests are algorithm-tagged---
\texttt{sha-256:9f86d0}\ldots, never a bare hex string---and a verifier handed
an untagged digest must refuse it rather than assume SHA-256. Assuming the
algorithm is exactly how a hash migration goes wrong, and it goes wrong
silently. And the claim set contains no floating-point numbers at all: \iat{}
is integer seconds, ordering comes from the step index. Number formatting is
the single hardest part of canonicalization, and the cheapest way to get it
right is to not have any numbers that need formatting.

Ten of the vectors are negative---each breaks exactly one rule, so that
rejecting it demonstrates one specific check exists. A vector rejected for the
\emph{wrong} reason is not a pass: a validator that chokes on the duplicate-key
vector because it could not parse the file has not shown that it detects
duplicate keys.

That duplicate-key vector is the one we would point at if we could only point
at one. It is a single file that a last-wins parser reads as \textsc{allow} and
a first-wins parser reads as \textsc{deny}. If the relying party and the
auditor make different choices there---and both choices are common, and neither
is wrong in any documentation---an operator can show each of them the answer it
wants and neither can tell. That is not a parsing preference. That is one
artifact with two meanings, which is the whole ballgame.

The reference implementation's own output is run through the validator in its
test suite, so the schema and the code cannot drift apart without something
going red. And because the CDDL is a real definition rather than an aspiration,
we implemented the CBOR rendering too and round-tripped every vector through
it. The two carry identical claim sets, and CBOR is \SI{48}{\percent} smaller,
because digests and signatures travel as bytes instead of as hexadecimal text.
That turns out to matter more than it sounds: as \S\ref{sec:pq} works out,
switching encoding saves more storage than post-quantum signatures cost.

\section{Writing It Down Precisely}
\label{sec:formal}

Words like ``the policy decides'' are fine until two people disagree about
what they mean, so here it is in symbols.

Let $\mathcal{A}$ be the actions an agent can propose, $\mathcal{I}$ the
identities it can act under, and $\mathcal{S}$ the states of the machine. The
steps taken so far form a path $p_t = \langle a_0, \dots, a_{t-1}\rangle \in
\mathcal{A}^{*}$. A policy is a function that looks at all four things and
gives one of four answers:
\begin{equation}
\Pi : \mathcal{A} \times \mathcal{I} \times \mathcal{S} \times \mathcal{A}^{*}
      \longrightarrow \mathcal{D},
\quad
\mathcal{D} = \{\textsc{allow},\ \textsc{deny},\
               \textsc{modify}(a'),\ \textsc{escalate}(h)\},
\end{equation}
where $a'$ is a cleaned-up version of the action and $h$ says which human to
ask. This function has to run inside the enclave, and it produces, for each
step, a proof
\begin{equation}
\pi_t = \mathrm{Sign}_{k_{\mathrm{AE}}}\!\left(
  H(\sigma_t^{\mathrm{snap}}) \,\|\, d_t \,\|\, t \,\|\, L_t \right),
\end{equation}
with $k_{\mathrm{AE}}$ a key nothing outside the enclave can use.

The fourth argument is the one that matters. Because $\Pi$ sees the path, the
read-then-send trick from Section~\ref{sec:path} is inside what the policy can
reason about, instead of being invisible to it.

And here is what we are careful \emph{not} to claim. We verify facts about
execution: which model ran, on what input, under whose authority, producing
what effect. We do not verify that the output was correct, or fair, or wise,
or that the boundaries somebody chose were the right boundaries. Those are
human judgments and they stay human judgments~\cite{iso42001}. It is easy to
let a system like this drift into implying more than it proves, and we would
rather draw the line ourselves than have someone else draw it for us later.

\subsection{This is information flow, and we should say so}
\label{sec:ifc}

Watching a path to catch read-then-send is not a new idea. It is information
flow control, which has been studied since the 1970s: Denning's lattice
model~\cite{denning1976,denning1977}, noninterference as
Goguen and Meseguer defined it~\cite{goguen1982}, and decades of systems that
carry labels around at the language and operating-system
level~\cite{myers1997,myers1999,sabelfeld2003,histar,flume}.

So rather than pretend we invented something, let us be exact about the
differences, because they are real:

\begin{itemize}[itemsep=2pt,topsep=2pt]
\item \textbf{How fine-grained.} Classical systems label variables and
processes inside a runtime they trust. We label \emph{actions at a boundary},
because the runtime here is a language model whose insides we specifically
decline to trust. Ours is coarser. It is also checkable from outside, which
theirs is not.
\item \textbf{What we are after.} Those systems \emph{enforce} a flow
property. We enforce a policy \emph{and hand you evidence} that we did, which
you can check without having watched. The evidence is the contribution; the
enforcement is an old idea done at a new place.
\item \textbf{What we promise.} No noninterference theorem. An agent whose
permitted actions are themselves harmful, or which leaks through a channel we
never intercept---timing, or the content of a permitted message---is outside
what we claim.
\item \textbf{What we inherit.} Label creep and the declassification
problem~\cite{sabelfeld2003}. A monitor that accumulates restrictions along a
path eventually refuses to let anyone do anything. The literature has been
saying this for fifty years. We did not believe it hard enough until we
measured it in Section~\ref{sec:utility}, where it cost us 42\% of legitimate
work.
\end{itemize}

\subsection{Keeping the cost from growing}
\label{sec:complexity}

Written as above, $\Pi$ takes the whole path as an argument, and the path
grows forever. If evaluating step $t$ costs something proportional to $t$,
then a long-running agent eventually stalls. That is not acceptable, so we
require policies to be expressible over a \emph{bounded summary}: a state
$\phi_t$ of size at most $B$, updated by folding in one step at a time,
\begin{equation}
\phi_t = \mathsf{upd}(\phi_{t-1}, a_{t-1}, d_{t-1}), \qquad
\phi_0 = \phi_{\mathrm{init}},
\end{equation}
so the implemented policy sees $\tilde{\Pi}(a, \iota, s, \phi_t)$ instead of
the whole history. The summary is doing the same job an IFC label does: it
carries the few facts the rules need---the highest classification read so far,
how much budget is left, whether anything has been declassified---and forgets
the rest.

Two things follow. The cost per step is $O(|\tilde{\Pi}| + B)$, flat no matter
how long the agent runs, and we commit $H(\phi_t)$ in the evidence so a
verifier can confirm the policy was evaluated against the right accumulated
state. But we also give something up: rules that need unbounded history
(``never contact an endpoint you did not see in the first $k$ steps'' for
arbitrary $k$) cannot be expressed exactly and must be approximated. That is a
real restriction. We would rather state it than quietly leave the design
unimplementable. Section~\ref{sec:eval} shows the flat cost is not just theory.

\section{What an Attacker Can Do}
\label{sec:security}

\subsection{Who we are up against}
We assume an attacker $\mathcal{A}$ who can run in polynomial time and who
(i)~owns the operating system, the container, the orchestration, and the
network---start, stop, reorder, replay, drop whatever it likes; (ii)~\emph{is
the operator}, and wants to look compliant. That second one is unusual and it
is the whole point: in an evidence system, the party writing the record is the
party being examined. Also (iii)~can put any text it likes in front of the
agent~\cite{greshake2023,injecagent}, and (iv)~can string together permitted
actions in any order.

What we assume it cannot do: \textbf{(A1)}~forge signatures or find hash
collisions; \textbf{(B1)}~get inside the enclave---and we should be honest
that this assumption has been dented in practice~\cite{costan2016,foreshadow};
\textbf{(C1)}~get the chip vendor to vouch for a lie, which is a
\emph{person} we are trusting, and exactly why vendor-rooted evidence alone
does not get above Tier~2; and \textbf{(D1)}~cut every verifier off from every
other verifier.

\subsection{Four things we want, and which ones we get}
\begin{description}[itemsep=2pt,leftmargin=1.4em]
\item[P1 (The verdict is real).] The decision in the token came from the
authorized rules applied to the description in the token. \emph{We get this}
under (A1), (B1), (C1): the measurement pins the code, the bundle hash pins
the rules, and the signature ties the verdict to the description.
\item[P2 (The evidence describes what happened).] \emph{We do not get this}
from attestation. This is snapshot substitution from
Section~\ref{sec:mediation}: without mediating the effect channel,
$\mathcal{A}$ produces a true statement about a fictional action. P1 holds
while P2 fails, which is the worst kind of failure---a perfect record of the
wrong thing.
\item[P3 (Nothing was changed or dropped).] \emph{We get this within a stretch
of history you actually see}: chaining catches edits, numbering catches holes.
\emph{We do not get it at the end}: if the operator simply stops showing you
the last few steps, nothing in the chain complains, because a chain cannot
tell you about a link you were never handed. Publishing the running total on a
bounded schedule fixes this, and shrinks the blind spot to the publishing
interval $\Delta$.
\item[P4 (Everyone sees the same history).] \emph{We do not get this} from
signing and publishing alone. The operator can keep two sets of books and show
each verifier a different one; both are internally perfect. This is the
split-view attack, well known in transparency
systems~\cite{rfc6962,ctgossip,chase2016}. The fix is that verifiers must
compare notes---gossip, a quorum of witnesses~\cite{chainiac}, or a ledger
that only admits one history.
\end{description}

Table~\ref{tab:threats} maps the threats to which property they attack.

\begin{table}[ht]
\centering\footnotesize
\begin{tabular}{@{}p{4.0cm}p{1.5cm}p{8.2cm}@{}}
\toprule
\textbf{Threat} & \textbf{Attacks} & \textbf{What stops it, and what is left over} \\
\midrule
Operator turns off enforcement or edits records & P1, P3 &
Rules and keys live in the enclave; changing the rules changes the
measurement and verifiers stop believing it. Left over: (B1), (C1). \\
Snapshot substitution & P2 &
Mediate the effect channel (Section~\ref{sec:mediation}). Left over: options 1
and 3 still trust the host not to fiddle with transport configuration. \\
Stop publishing recent history & P3 &
Bounded publishing interval $\Delta$; a missing publication is itself an
alarm. Left over: whatever happened in the last $\Delta$. \\
Two sets of books & P4 &
Gossip, witnesses, or consensus. Left over: witnesses who collude. \\
Injected instructions from the world & P1 (indirectly) &
The enclave applies fixed rules and does not read persuasion, so an injected
instruction cannot authorize an out-of-scope action. Left over: harmful
actions that are inside scope---a safety problem, not a verifiability
one~\cite{greshake2023,toolemu}. \\
Stringing permitted actions together & P1 &
The policy sees the path (Section~\ref{sec:formal}). Left over: the cost in
refused legitimate work, measured in Section~\ref{sec:utility}. \\
Sub-agents used as cutouts & P1--P4 &
Constraints are inherited; sub-agent evidence checked by several
verifiers~\cite{multiverifier}. \\
Breaking the enclave & P1--P4 &
Out of scope by (B1)---disclosed, not solved, and bounded by re-attestation
and key rotation. \\
\bottomrule
\end{tabular}
\caption{Threats against properties. P2 and P4 are the ones you do
\emph{not} get for free, which is why the standard has extra requirements for
them.}
\label{tab:threats}
\end{table}

\subsection{Proving the two hard ones}
\label{sec:proofs}

P1 and P3 are ordinary signature-and-chain arguments. P2 and P4 are the
interesting ones, so we state them properly.

\begin{definition}[Evidenced execution]
\label{def:evidenced}
Fix an agent $\iota$ and a relying party $R$. An effect $e$ carried out at $R$
is \emph{evidenced} if some token $\tau$ that a verifier accepts commits to a
description digest $H(\sigma^{\mathrm{snap}})$, a target $r$, and the verdict
\textsc{allow}, where $\sigma^{\mathrm{snap}}$ describes exactly $e$ and $r$
is where $e$ happened.
\end{definition}

\begin{definition}[The execution-fidelity game $\mathsf{Exp}^{\mathrm{fid}}$]
$\mathcal{A}$ has the powers listed above and may talk to the enclave
$\mathcal{E}$ as much as it likes, collecting tokens and tickets. It wins if a
compliant relying party $R$ carries out an effect $e$ that is not evidenced in
the sense of Definition~\ref{def:evidenced}.
\end{definition}

In plain words: the attacker wins if it can get something done that the
paperwork does not describe.

\begin{theorem}[Tickets give you execution fidelity]
\label{thm:fidelity}
Suppose for each \textsc{allow} the enclave issues
$c = \mathrm{Sign}_{k_{\mathcal{E}}}(H(\sigma^{\mathrm{snap}}) \,\|\, r
\,\|\, \mathrm{nonce})$, and a compliant $R$ acts only when the request
carries a $c$ that (i)~verifies under an enclave key whose measurement matches
the published value, (ii)~names $R$'s own resource, (iii)~carries an unused
nonce, and (iv)~matches the request itself against
$H(\sigma^{\mathrm{snap}})$. Then for any polynomial-time $\mathcal{A}$,
\[
\Pr[\mathcal{A} \text{ wins } \mathsf{Exp}^{\mathrm{fid}}]
\;\le\; \mathbf{Adv}^{\mathrm{EUF\text{-}CMA}}_{\mathrm{Sig}}(\mathcal{A})
\;+\; \mathbf{Adv}^{\mathrm{coll}}_{H}(\mathcal{A}),
\]
under (A1), (B1), (C1).
\end{theorem}

\begin{proof}[Proof sketch]
Say $\mathcal{A}$ wins: $R$ did $e$, and no token describes $e$ at $r$. By
check~(i) the ticket verified under the enclave key, which by (B1) never
leaves and by (C1) is tied to the attested measurement. So either
$\mathcal{A}$ forged a signature---that case is bounded by
$\mathbf{Adv}^{\mathrm{EUF\text{-}CMA}}$---or the enclave really issued the
ticket. If it really issued it, then by check~(iv) the request that was
carried out hashes to the digest inside the ticket. If the description had
been of something other than $e$, two different descriptions would share a
digest, which is bounded by $\mathbf{Adv}^{\mathrm{coll}}_{H}$. Otherwise the
description \emph{is} of $e$; by check~(ii) the target matches; and the
enclave only issues tickets alongside an \textsc{allow} token it also emits.
So $e$ is evidenced---contradiction. Check~(iii) stops a ticket being reused
for a second effect. \end{proof}

\begin{proposition}[Without the far-end check you lose immediately]
\label{prop:necessity}
If nothing at the point of effect ties the request to
$H(\sigma^{\mathrm{snap}})$, then an attacker owning the host wins with
probability 1.
\end{proposition}
\begin{proof}
Hand the enclave a description of something harmless, take the genuine
\textsc{allow}, then use the host's own credentials to do something else.
No check anywhere is violated: the token is real and the enclave's statement
is true. The effect is unevidenced, so the attacker wins every time. \end{proof}

We want to be blunt about what Proposition~\ref{prop:necessity} means, because
it is the thing we would most like a reader to take away. It is not a corner
case or a hardening tip. It says that an evidence architecture without effect
mediation is broken at probability one while every signature verifies
perfectly. That is why the requirement is a MUST.

\begin{definition}[Equivocation]
$\mathcal{A}$ equivocates if two verifiers accept histories for the same agent
where neither is a prefix of the other.
\end{definition}

\begin{theorem}[Two sets of books can be caught, and pinned on someone]
\label{thm:equiv}
Suppose each running total at index $i$ is signed by the enclave, and two
verifiers exchange (total, index) pairs---directly, through a witness quorum
with at least one honest member, or via a single-history ledger. If
$\mathcal{A}$ equivocates, then under (A1) and (D1) the honest parties end up
jointly holding two validly signed totals at the same index with different
values. That is a non-repudiable proof of equivocation, attributable to the
enclave key. Detection probability is 1.
\end{theorem}

\begin{proof}[Proof sketch]
Two histories that are not prefixes of one another differ at some earliest
index $i$. Both verifiers hold signed totals at $i$; when they compare, the
values differ. Forging either reduces to signature forgery, ruled out by (A1).
And the pair stands on its own---you do not have to believe either verifier's
story about it. \end{proof}

Notice what Theorem~\ref{thm:equiv} does \emph{not} say. Gossip does not
prevent an operator from keeping two sets of books. It makes it impossible to
keep them secret, and it produces a receipt with a name on it. That is the
same guarantee Certificate Transparency provides~\cite{rfc6962,chase2016}, and
it is worth being precise about, because ``we use gossip'' sounds like
prevention and is not.

\subsection{Building Tier 4}
\label{sec:tier4}
Tier~4 says the system cannot act unless verification succeeds. That has to be
a construction, not a promise, and it falls out of the ticket scheme. If the
far end refuses any request without a valid ticket, then: when the enclave is
not running, or its measurement does not match, no ticket exists, so nothing
happens. The system fails closed \emph{because of how it is built}, not
because somebody left a flag set. And the decision to stop is not the
operator's to make---which is the difference between this and a ``halt''
variable in a process the operator controls and can patch out. What is left
over is the far end's own enforcement and the chip vendor (C1), both of which
must appear in the disclosure.

This also settles a tension people notice: a TEE attestation on its own is
Tier~2, because (C1) is a person you are trusting. You get to Tier~3 by adding
independent anchoring, so nobody has to take the vendor's word about which
code ran, and to Tier~4 by making valid evidence a precondition of acting. The
tiers grade the shape of the trust, not the brand of the cryptography.

\subsection{Will anyone actually do this?}
\label{sec:deploy}

Theorem~\ref{thm:fidelity} asks the far end to check tickets, and it is fair
to ask whether that is realistic. Three things.

\emph{The mechanism is old.} Action-bound tickets are capabilities, which have
been around since the 1960s~\cite{dennis1966,levy1984}, and macaroons already
give you exactly this shape---bearer tokens with cryptographic caveats,
checkable by the service without calling anyone~\cite{macaroons}, over
ordinary HTTP headers. A check is a signature verification and a comparison,
not a new industry.

\emph{You can start alone.} Options 1 and 3 need nobody's cooperation: run
egress from inside the enclave, or force it through an attested proxy. Inside
one company, the API gateway is already under your control, so you can turn on
ticket checking there first. Partners come next by contract---which is how SOC
2 requirements already propagate---and standardization last, which is what the
attestation-token profile work is for.

\emph{Some endpoints will never do it.} Arbitrary third-party services and
legacy systems can only be covered by options 1 or 3, whose guarantee is
weaker. We are not claiming a universal answer. We are claiming a property,
three ways to get it with different amounts left over, and a rule that you
have to say which one you used---so the person relying on your evidence can
tell the difference.

\section{How Much of This Do Existing Rules Already Ask For?}
\label{sec:gap}

Before claiming a gap exists, you should check whether somebody already filled
it. So we took all 127 requirements and, one at a time, asked of eight
existing frameworks: does this framework already require this?

\subsection{How we did it}
Using the coding method HAARF worked out for healthcare
agents~\cite{haarf2026}, each requirement gets one of three marks against each
framework. \emph{Exact}: they ask for the same thing, at the same depth.
\emph{Partial}: they ask about the same topic, but without demanding evidence
somebody outside the company can check, or not in as much detail.
\emph{None}: they do not address it. Coverage is (Exact + Partial) over 127.
The sheet, the rubric, and the script are public, so you can re-derive the
numbers or re-code a framework you think we got wrong.

\subsection{Checking the coding with a different model}
\label{sec:crossmodel}

The requirements and the coding sheet were produced by the same party. That is
a structural weakness in any coverage mapping: whoever decides that an external
framework does not already cover a requirement is the party that wrote the
requirement. The two-coder study the rubric calls for is the answer to it, and
it has not been run.

What can be run cheaply is a cross-model audit. A model from a different vendor
reviews every row---one independent pass per framework, without access to the
original reasoning, instructed to look for miscodings in both directions and to
name a specific clause for every proposal. The procedure matters more than the
idea, because the cheap version of it is worthless.

The failure mode is fabrication. A model asked for a clause citation will
produce a plausible-looking one, and a fabricated citation is worse than no
citation because it survives casual review. Three controls address it.

\paragraph{An acceptance rule fixed in advance.} A proposal is applied only if
it can be verified \emph{without} trusting the reviewer about a framework's
contents; everything else is recorded rather than applied. Here that admitted
31 of 192. Twenty-four were Exact ratings on the conformance-statement
requirements, settled by reading the requirement text: those describe an
artifact this standard defines, so nothing external can be equivalent in scope
and intent to it, and no knowledge of SOC 2 is needed to establish that. Seven
cited MITRE ATLAS mitigation identifiers confirmed to exist. The remaining
161---including all 91 proposed \textsc{partial}-to-\textsc{none}
downgrades---wait for a human with the source document open.

\paragraph{A check on the direction of the corrections.} Applying every
proposal would have removed about 64 covered rows and taken 9 to 15 points off
most frameworks: the direction that widens the gap this standard exists to
fill. A review that finds errors in only one direction is measuring the
reviewer rather than the sheet. This one proposed 37 changes that \emph{raised}
external coverage, which is what makes its downgrades worth reading.

\paragraph{A factual spot-check.} One claim was checkable against a published
source---that the MITRE ATLAS coding drew on a superseded mitigation set.
Current ATLAS includes AI Bill of Materials (AML.M0023), AI Telemetry Logging
(M0024), Maintain AI Dataset Provenance (M0025), and agent-specific permission
mitigations (M0026--M0028). Coding them moved ATLAS coverage from 25\% to
30\%---upward, against the direction that would flatter this standard.

\paragraph{What it returned.} 192 proposed changes and 96 citations that do not
name a checkable clause. Three results are structural rather than row-by-row.
Most citations in the sheet point at a chapter heading rather than a
requirement---for two frameworks, essentially all of them---so a reader cannot
look up the claim being made. Six hundred and ninety-two of the thousand rows
carry a rationale repeated word-for-word on four or more rows, the mechanical
signature of coding by block rather than by requirement. And every one of the
61 Exact ratings was challenged, which is consistent with this paper's own
thesis: if no existing framework requires operator-independent evidence, the
Exact column should be close to empty.

The pass cost a few minutes of compute. A cross-model audit with a stated
acceptance rule is worth running for any coverage mapping, coding study, or
requirements traceability matrix where the author and the coder are the same
party---not because a second model is an authority, but because it is an
independent generator of specific, checkable objections. It produces no
$\kappa$, it does not replace human coders, and it should never be applied
wholesale. The prompt, the output schema, the 192 proposals with per-item
confidence, and the record of what was applied and what was refused are in the
repository, so the pass can be repeated or contradicted.

\subsection{What came out}
Figure~\ref{fig:coverage} and Table~\ref{tab:coverage} have the results, and
two patterns are worth staring at.

The first is that the Partial column is enormous and the Exact column is
nearly empty. \emph{Five of the eight frameworks have zero exact matches}, and
the NIST risk framework---which scores 60\%, second highest overall---is one of
them. Read that again, because it is the whole story of this paper in one
number. These frameworks ask for nearly every control we ask for. What they do
not ask for, essentially anywhere, is machine-made evidence that the control
actually held, checkable by someone who does not trust you.

The second is where the None column bunches up: in the chapters about evidence
properties, grading, and disclosure. No framework we coded grades evidence by
how independently it can be checked, or requires you to write down what you
are still asking people to take on faith.

\begin{figure}[t]
\centering
\includegraphics[width=0.9\linewidth]{figures/mapping-coverage.pdf}
\caption{How much of \poc{} eight frameworks already require. Dark: exact
matches. Light: partial. The rest is the gap.}
\label{fig:coverage}
\end{figure}

\begin{table}[ht]
\centering\small
\begin{tabular}{@{}lcccc@{}}
\toprule
\textbf{Framework} & \textbf{Exact} & \textbf{Partial} & \textbf{None} &
\textbf{Coverage} \\
\midrule
OWASP AISVS~\cite{owaspaisvs} & 16 & 63 & 48 & 62\% \\
NIST AI RMF~\cite{nistairmf} & 0 & 75 & 52 & 59\% \\
ISO/IEC 42001~\cite{iso42001} & 0 & 72 & 55 & 57\% \\
SOC 2~\cite{soc2tsc} & 0 & 68 & 59 & 54\% \\
EU AI Act~\cite{euaiact} & 0 & 65 & 62 & 51\% \\
CSA AARM~\cite{csaaarm} & 13 & 47 & 67 & 47\% \\
Zero Trust (SP 800-207)~\cite{sp800207} & 8 & 46 & 73 & 43\% \\
MITRE ATLAS~\cite{mitreatlas} & 0 & 38 & 89 & 30\% \\
\bottomrule
\end{tabular}
\caption{Coverage of all 127 requirements, coded individually by one coder.
The thin Exact column is the finding: frameworks want the controls, not
operator-independent evidence that the controls held.}
\label{tab:coverage}
\end{table}

\subsection{Why you should not lean too hard on these numbers}
\label{sec:gaplimits}

Now, a warning, and we mean this seriously. It is very easy to fool yourself
with a table like that, and the person most likely to be fooled is us.

\emph{The measurement points the way we wanted it to.} We asked how much of
\emph{our} list other people cover. Invent any sufficiently novel list and
everyone else scores badly on it. So this table says we are working somewhere
different. It does not say we are working somewhere \emph{useful}. That is a
separate argument and these numbers do not make it.

\emph{One person coded it.} There is no measure of agreement at all. The published rubric specifies what a
real study looks like: two independent coders per framework, Cohen's
$\kappa$ reported~\cite{cohen1960}, at least one coder who did not write the
requirements, and disagreements resolved by a third. That study is planned, and
when it is done it replaces Table~\ref{tab:coverage}.

The cross-model audit in \S\ref{sec:crossmodel} is a partial substitute and
not a replacement: it yields no measure of agreement, and its findings on
citation quality and block-level coding apply to the sheet behind
Table~\ref{tab:coverage} as it stands.

\emph{The grain is coarser than it looks.} Coding is per requirement, but most
sections are still coded uniformly across the requirements they contain, so the
apparent resolution of the sheet exceeds its real resolution. The published
crosswalks now summarize each section over the requirements beneath it and name
the ones a framework does not reach.

\emph{We only looked one way.} We did not measure what those frameworks
require that \emph{we} miss---and there is plenty: organizational governance
in the NIST framework, management-system obligations in ISO/IEC~42001,
availability and processing integrity in SOC~2. We are deliberately narrower
than all of them. An honest comparison shows both directions, and ours only
shows one so far.

\subsection{Compared with the closest things}
Table~\ref{tab:compare} compares us against the systems we are most often
confused with. The comparison is on \emph{specified} properties---what each
approach requires an implementation to do---because that is the only fair
comparison available. ACS and AARM~\cite{acs2026,acsms2026,csaaarm} specify
interception and enforcement, then leave the record with the operator.
Transparency logs~\cite{rfc6962,chainiac} have the non-equivocation machinery
but know nothing about agents. Confidential-computing pipelines~\cite{eqty2026}
attest computation but have no notion of an agent's authority or delegation.
What we are contributing is the combination and the grading rule, not any one
of the parts.

\begin{table}[ht]
\centering\footnotesize
\begin{tabular}{@{}p{2.35cm}p{2.2cm}p{2.2cm}p{2.4cm}p{2.4cm}p{2.4cm}@{}}
\toprule
\textbf{Property} & \textbf{Prompts / guardrails} & \textbf{Static RBAC} &
\textbf{ACS / AARM} & \textbf{Transparency logs} & \textbf{\poc{}} \\
\midrule
Every action goes through it & No & At the API layer & Yes & N/A & Yes \\
Decisions consider the path & No & No & Partly & N/A & Yes \\
Record integrity & None & Operator's logs & Operator's receipts &
Chained, gossiped & Chained + anchored \\
A stranger can check it & No & No & With operator access & Yes &
Yes (Tier 3+) \\
Everyone sees one history & No & No & Unspecified & Yes (with gossip) &
Required at Tier 3+ \\
Evidence describes what happened & No & No & Unspecified & N/A &
Required (mediation) \\
Says what you still trust & No & No & No & Implicitly & Standardized \\
\bottomrule
\end{tabular}
\caption{What each approach \emph{requires}. We are a specification; no claim
is made here about deployed performance.}
\label{tab:compare}
\end{table}

\section{Does It Actually Work?}
\label{sec:eval}

Talking about a design is easy. We built it and measured it, because that is
the only way to find out whether the arguments above survive contact with a
computer. The code is open source in the specification
repository~\cite{pocspec}, at \url{https://github.com/AAI-Society/ov-poc-standard}.

\paragraph{What is real, and what is pretend.} Real: the signing and
verification, the hashing and chaining, the canonical descriptions, the
path-aware policy over a bounded summary, the tickets and the far-end checks,
the anchoring, the gossip, and the fail-closed behavior when the evidence
store goes away. Pretend, \emph{in this section}: the hardware. The
measurements below use an enclave that is an in-process object whose key is
never handed out and whose ``measurement'' is a digest of the policy code. That
reproduces the structure the theorems are about, but not hardware isolation.
\textbf{So the timings below leave out the cost of crossing into a real
enclave, and the real number is bigger.} We would rather report a lower bound we
measured than a total we invented---and we have since measured the difference
rather than leaving it as a caveat: \S\ref{sec:tdx} runs the same pipeline
inside a real Intel TDX trust domain and puts the trust-domain overhead at
6.2\%. Everything below: Apple M2 Max, Python 3.11, one core, 5{,}000
iterations after 200 warm-ups.

\subsection{The attacks}
We wrote eleven attacks and ran each one twice: once against a system missing
the relevant requirement, once against a system with it
(Table~\ref{tab:attacks}). All eleven work when the requirement is missing and
all eleven are refused or caught when it is there. The last three were not
planned. A9 fell out of building the Merkle experiment in \S\ref{sec:merkle};
A10 and A11 came from an independent review of the implementation, and both
were live defects in it.

\begin{table}[ht]
\centering\footnotesize
\begin{tabular}{@{}p{0.5cm}p{3.3cm}p{4.6cm}p{5.2cm}@{}}
\toprule
& \textbf{Attack} & \textbf{Without the requirement} & \textbf{With it} \\
\midrule
A1 & Snapshot substitution & Works first try: the record describes a harmless
action while a malicious one runs & Refused: request does not match the
evidenced description \\
A2 & Editing the log & Works against unsigned operator logs &
Caught: bad signature at record 2 \\
A3 & Dropping a step & Works: a dropped step looks like a step that never
happened & Caught: gap at step 3 \\
A4 & Hiding recent history & Works: the shortened history is internally
perfect & Caught: published total covers 8 steps, 5 shown \\
A5 & Two sets of books & Works: each verifier's copy is internally perfect &
Caught: different totals at the same index \\
A6 & Read then send & Works: each action is inside the grant &
Blocked: egress after a confidential read \\
A7 & Reusing a ticket & Works without single-use tickets &
Refused: nonce already used \\
A8 & Evidence pipeline down & Works: the system keeps acting while logging is
broken & Refused: fails closed, failure recorded \\
A9 & Rewriting history, by someone who holds the key & Works: the rewritten
history replays perfectly---every signature valid, every link correct, because
the operator recomputed them all & Caught: the head presented does not match the
total published \emph{before} the rewrite \\
A10 & Substitution past a check that is switched on but unconfigured & Works:
the endpoint verifies signature, measurement, resource and nonce, and still runs
a different request & Refused: the ticket carries a digest of the action the
endpoint can recompute alone, and the check refuses when it cannot \\
A11 & A valid ticket shown beside the wrong record & Works: both objects verify
on their own, so the reconciliation balances & Caught: the ticket names the step
and record it was issued with \\
\bottomrule
\end{tabular}
\caption{The attack harness. A1 is Proposition~\ref{prop:necessity} in the
flesh: without the far-end check, the attacker wins on the first attempt. A10 is
the same proposition in the form it takes in practice---the check present,
enabled, and doing nothing.}
\label{tab:attacks}
\end{table}

\subsection{What it costs}
Table~\ref{tab:latency} breaks a single intercepted action into pieces. The
signing is 87\% of it. The policy evaluation---the part everyone worries
about---is \SI{0.37}{\micro\second}, which is nothing. The total,
\SI{201}{\micro\second}, is 1.3\% of our \SI{15}{\milli\second} budget. Put
the other way: there is 74 times more room than we are using, and that room is
what a real enclave transition would eat into. One core does 4{,}963 actions a
second.

(The total is bigger than the sum of the parts because a full step signs twice
---once for the record, once for the ticket---and serializes the token.)

\begin{table}[ht]
\centering\small
\begin{tabular}{@{}lrrr@{}}
\toprule
\textbf{Piece} & \textbf{Mean} & \textbf{p95} & \textbf{p99} \\
\midrule
Building and hashing the description & \SI{3.72}{\micro\second} & \SI{4.21}{\micro\second} & \SI{8.00}{\micro\second} \\
Deciding (path-aware) & \SI{0.37}{\micro\second} & \SI{0.46}{\micro\second} & \SI{1.12}{\micro\second} \\
Signing (Ed25519) & \SI{84.25}{\micro\second} & \SI{97.96}{\micro\second} & \SI{110.12}{\micro\second} \\
Adding the link, and the tree root & \SI{0.73}{\micro\second} & \SI{0.92}{\micro\second} & \SI{3.71}{\micro\second} \\
\midrule
\textbf{One whole action} & \textbf{\SI{201.47}{\micro\second}} & \textbf{\SI{230.58}{\micro\second}} & \textbf{\SI{254.75}{\micro\second}} \\
\bottomrule
\end{tabular}
\caption{Cost of one intercepted action, not counting enclave transitions.}
\label{tab:latency}
\end{table}

\subsection{Does it stay cheap as the agent runs longer?}
Section~\ref{sec:complexity} claimed the bounded summary makes the per-step
cost independent of how long the agent has been going. That is the kind of
claim that is easy to make and easy to be wrong about, so we measured it
(Table~\ref{tab:scaling}). From 10 steps to 50{,}000---a factor of five thousand---the
cost per step does not move: it sits at \SI{0.21}{\micro\second} throughout
(Figure~\ref{fig:scaling}). The naive version, which re-reads the whole history
every time, grows exactly as you would expect and would blow the entire
\SI{15}{\milli\second} budget somewhere around 22{,}000 steps.

So this is not bookkeeping. It is the difference between a monitor that
survives a long-running agent and one that gets slower until somebody turns it
off.

\begin{figure}[t]
\centering
\includegraphics[width=0.88\linewidth]{figures/chart-scaling.pdf}
\caption{Deciding one action, at depth, on log axes. The bounded summary is a
flat line across four orders of magnitude; re-reading the whole path is a
straight diagonal that crosses the \SI{15}{\milli\second} budget at roughly
22{,}000 steps.}
\label{fig:scaling}
\end{figure}

\begin{table}[ht]
\centering\small
\begin{tabular}{@{}rrr@{}}
\toprule
\textbf{Steps so far} & \textbf{Bounded summary} & \textbf{Re-reading everything} \\
\midrule
10 & \SI{0.208}{\micro\second} & \SI{5.79}{\micro\second} \\
100 & \SI{0.208}{\micro\second} & \SI{64.92}{\micro\second} \\
1{,}000 & \SI{0.208}{\micro\second} & \SI{663.88}{\micro\second} \\
5{,}000 & \SI{0.209}{\micro\second} & \SI{3379.17}{\micro\second} \\
50{,}000 & \SI{0.208}{\micro\second} & \SI{34368.50}{\micro\second} \\
\bottomrule
\end{tabular}
\caption{Cost of deciding one action, at depth. Flat versus linear.}
\label{tab:scaling}
\end{table}

Checking, from outside, costs \SI{190}{\micro\second} per record and grows
linearly: a stranger can verify a 5{,}000-step history in under a second using
nothing but published keys and reference values. The important word there is
\emph{stranger}. No cooperation from the operator is required at any point.

\subsection{The measurement we did not expect}
\label{sec:utility}

Everything above says the machinery is cheap. This section says the policy is
expensive, in a way no amount of faster code will fix.

We ran 2{,}000 randomized workflows---70\% ordinary work, 30\% trying the
read-then-send trick---past the path-aware monitor and counted how often it
refused work it should have allowed. With no way to clear an accumulated
label, it refused \textbf{42\%} of perfectly legitimate work. Not a bug: we
told it ``don't send after reading something sensitive,'' and it enforced
exactly that, forever, because nothing ever cleared the flag.

The standard remedy, from fifty years of information-flow
work~\cite{sabelfeld2003}, is a \emph{declassification point}: a place where
data is redacted and the accumulated label is reset. Adding one took the false
rejections to zero.

\subsection{Where we fooled ourselves, and the fix}
\label{sec:frontier}

That result was too good, and it was too good because we cheated without
noticing. In that first experiment only the \emph{benign} workflows were
allowed to declassify. A real attacker would of course use any declassification
point available to it---that is precisely what makes declassification hard, and
we had quietly assumed the problem away.

So we ran it again properly. This time a declassification point is available to
whoever reaches it, and the monitor---which sees only that a redaction step
occurred, and cannot tell a genuine redaction from a sham---treats it the same
either way. We swept the \emph{coverage}: what fraction of workflows have a
declassification point before egress. Figure~\ref{fig:frontier} has the result,
and it is much less flattering.

\begin{figure}[t]
\centering
\includegraphics[width=0.88\linewidth]{figures/chart-frontier.pdf}
\caption{The declassification frontier. Dashed: declassification the monitor
cannot verify---every point of relief from false rejections costs about 1.6
points of detection, and there is no knee to aim for. Solid: declassification
the monitor can verify, where both quantities improve together.}
\label{fig:frontier}
\end{figure}

The dashed line is a straight trade, not a bargain. Going from no
declassification to full coverage takes false rejections from 44.8\% to zero
and takes detection from 71.7\% to zero along with it---\textbf{1.6 points of
detection lost for every point of relief}. There is no knee. There is no good
place to sit. An unverified declassification point is not a solution to label
creep; it is a dial between two different ways of failing.

The fix is not to tune the dial. It is to make declassification something the
monitor can check: only a trusted redaction tool clears the label, and it
really removes the content, so an attacker routing through it gains nothing
because the data it wanted is gone. That is the solid line, and it does what
you want in both directions at once: at full coverage, false rejections reach
zero \emph{and} detection reaches 100\%, because a redaction that genuinely
redacts defeats the exfiltration it was supposed to launder.

This changes a requirement, not a setting. \textbf{A deployment that mandates
path-aware authorization must specify not only where declassification happens
but how the monitor verifies it.} Declassification by annotation---a workflow
asserting ``this is fine now''---measurably destroys the detection the whole
mechanism exists to provide.

We also checked whether the size of the bounded summary mattered, sweeping the
label bound $B$ from 1 to 16. It made no difference at all on these workloads:
the binding constraint is the declassification policy, not the size of the
state we carry. That is a null result and we report it as one.

\subsection{Reasons to doubt these numbers}
The enclave is pretend here, so a real deployment costs more than
\SI{201}{\micro\second}---by 6.2\% for the trust domain itself, which
\S\ref{sec:tdx} measures on real Intel TDX, plus \SI{39.5}{\milli\second}
for each hardware quote, which is the cost that actually matters and is not in
any number in this section. Everything runs in one process, so no network or IPC
costs are included. The workflows are synthetic: the shape of the utility
result is solid---label creep is real and declassification fixes it---but
42.2\% is a property of our workload model and yours will differ. It is
Python, so the absolute numbers are conservative; the breakdown between
components is the part that will hold up. And we did not implement the
privacy-preserving evidence options at all.

Four more, specific to the newer work. Everything is measured on \emph{one}
core with \emph{one} log, and a hash chain is inherently serial---each link
depends on the one before---so a single global log will cap throughput no
matter how many cores you have. Per-agent logs parallelize but make cross-agent
ordering, and therefore delegation, harder to reason about; the standard does
not yet say which you must build, and it should. Our tree keeps every leaf in
memory so it can generate proofs, which is fine for a log that stores the
records anyway but is not a storage design. The CBOR claim keys are a
provisional private-use allocation pending IANA registration, so the round-trip
property is what we are asserting, not the integers. And our reference
serializer is not quite RFC 8785: it orders keys by code point rather than by
UTF-16 code unit, which is identical for the claim set, whose keys are all
ASCII, and can differ for an application payload with a key outside the Basic
Multilingual Plane. We publish a test vector that exhibits exactly that
difference rather than describing it and hoping.

\subsection{Signatures that have to outlive the attacker}
\label{sec:pq}

There is one more thing worth measuring, and it comes from taking the purpose
of the system seriously. Evidence is not a session key. A record made today
may be looked at in a dispute, an audit, or a courtroom ten or twenty years
from now, and it is worth exactly what its signature is worth \emph{at the
moment somebody looks}. So a signature scheme with an expiry date does not
merely expire. It reaches back and destroys the value of evidence about things
that happened years earlier. For most systems post-quantum migration is about
protecting future traffic; for an evidence system it is about protecting the
past.

The good news is that most of the machinery does not care. The hash chain, the
sequence numbers, and the anchoring are already post-quantum: the best known
quantum attack on preimages is Grover's, which buys a square root, leaving
SHA-256 at about 128 bits. So the structure survives. What has to migrate is
the signatures---on tokens, on capabilities, and on published roots.

We measured what that swap costs (Table~\ref{tab:pq}), using NIST's ML-DSA
(FIPS 204) against Ed25519.

\begin{table}[ht]
\centering\small
\begin{tabular}{@{}lrrrr@{}}
\toprule
\textbf{Scheme} & \textbf{Signature} & \textbf{Public key} &
\textbf{Record} & \textbf{Per million actions} \\
\midrule
Ed25519 (classical) & 64 B & 32 B & 1{,}336 B & \SI{1.34}{\giga\byte} \\
ML-DSA-44 & 2{,}420 B & 1{,}312 B & 10{,}760 B & \SI{10.76}{\giga\byte} \\
ML-DSA-65 & 3{,}309 B & 1{,}952 B & 14{,}316 B & \SI{14.32}{\giga\byte} \\
Hybrid Ed25519 + ML-DSA-44 & 2{,}484 B & 1{,}344 B & 11{,}016 B & \SI{11.02}{\giga\byte} \\
\bottomrule
\end{tabular}
\caption{Cost of post-quantum evidence signatures. Records carry two
signatures (token and capability), hex-encoded. Sizes are
implementation-independent; see the caveat below about timing.}
\label{tab:pq}
\end{table}

The headline is size, not speed. A signature goes from 64 bytes to 2{,}420---%
a factor of 38---and because an evidence chain stores one per action, the
whole record grows about eightfold, from roughly \SI{1.3}{\giga\byte} to
\SI{10.8}{\giga\byte} per million actions. Before you accept that bill, note
that switching from the JSON rendering to the CBOR one halves every figure in
the table (\S\ref{sec:interop}), because a signature stored as hexadecimal
text costs twice what the same signature costs as bytes. The encoding choice is
worth more than the algorithm choice costs. That is the real bill, and it is a
storage and bandwidth conversation, not a cryptography one.

\paragraph{A caveat we want to be loud about.} Our Ed25519 is a C-backed
library and our ML-DSA is a pure-Python reference implementation, so the
timings are not comparable and we are deliberately not putting them in the
table. Pure-Python ML-DSA signing measured tens of milliseconds, which would
blow the whole latency budget---but that number says something about our
Python, not about ML-DSA. Optimized implementations are competitive with
Ed25519, and verification in particular is fast. Anyone who cites a timing
comparison from this section is citing an artifact. The size numbers are the
ones that transfer.

Figure~\ref{fig:retention} puts that in units anyone can budget for. A fleet
of a thousand agents taking a hundred actions each per day, kept for the seven
years a regulator might ask about, comes to \SI{0.24}{\tera\byte} classically
and \SI{2.65}{\tera\byte} post-quantum. The multiplier is 38 on the
signature and about eleven on the whole record; the absolute number is a couple
of terabytes, which is to say the frightening-sounding factor turns out to be a
rounding error on a storage invoice. We would rather have measured that than
argued about it.

\begin{figure}[t]
\centering
\includegraphics[width=0.88\linewidth]{figures/chart-retention.pdf}
\caption{What keeping the evidence costs, for a realistic fleet. Post-quantum
signatures multiply the bill by about eleven; the bill is small enough that it
does not matter.}
\label{fig:retention}
\end{figure}

\paragraph{What we changed in the standard.} Two requirements. First, every
record and capability must name the algorithm that signed it, and a claim must
declare a migration path---so a verifier knows what to check and an operator
can change schemes without orphaning evidence already published. Second, if
the declared retention period runs past the period the signature scheme is
expected to survive, then the implementation either signs post-quantum or
hybrid, or re-signs and re-anchors retained evidence before the projection
lapses. Hybrid is the sensible default during transition: you pay 2{,}484
bytes instead of 2{,}420 and you are covered whichever way the next decade
goes.

\subsection{Choosing the anchoring interval, by measurement}
\label{sec:deltameasured}

Earlier we said the publishing interval $\Delta$ bounds how much recent history
an operator can hide, and we offered a table of suggestions. Suggestions are
unsatisfying in a paper that measures everything else, so we measured this too.

Publishing one root costs \SI{85.7}{\micro\second}---one signature and an
append. Our gateway sustains about 4{,}989 actions per second on one core.
Those two numbers determine everything (Figure~\ref{fig:anchor}): the blind
window is the action rate times $\Delta$, and the overhead is the publication
cost divided by $\Delta$.

\begin{figure}[t]
\centering
\includegraphics[width=0.88\linewidth]{figures/chart-anchor.pdf}
\caption{The anchoring tradeoff, measured. Exposure and cost are exactly
inverse, so the only question is which one you would rather spend.}
\label{fig:anchor}
\end{figure}

The numbers are friendlier than we expected. Publishing every second leaves at
most about 5{,}000 actions inside the blind window and costs
$8.6 \times 10^{-5}$ of wall-clock time---under a hundredth of a percent.
Publishing every millisecond shrinks the window to five actions and still costs
under 9\% of one core. Publishing hourly is essentially free and leaves 18
million actions hideable, which for anything irreversible is absurd.

The guidance becomes arithmetic rather than taste: \emph{decide how many
actions you could tolerate an adversary concealing, divide by your action rate,
and that is your $\Delta$.} For irreversible actions the tolerable number is
near zero, which is why we recommend publishing per action there and paying the
round trip.

\subsection{Getting rid of the signing cost}
\label{sec:batch}

Signing is 87\% of a step, which invites an obvious question: why sign every
action instead of one root over a batch of them? We measured that too
(Figure~\ref{fig:batch}).

\begin{figure}[t]
\centering
\includegraphics[width=0.88\linewidth]{figures/chart-batch.pdf}
\caption{Batch signing. Throughput rises almost linearly with batch size while
the window in which an action is not yet covered by a signature stays in the
microseconds.}
\label{fig:batch}
\end{figure}

Batching 128 actions under one signature makes each action \textbf{64 times
cheaper}---from \SI{85.7}{\micro\second} to \SI{1.3}{\micro\second}, and from
11{,}663 to 745{,}573 actions per second on one core. The price is that an
action is not covered by a signature until its batch closes, which at that size
is at most \SI{170}{\micro\second}.

That is the same shape of tradeoff as $\Delta$, one level down: a small window
in which the record exists but is not yet cryptographically fixed. If the
machine dies inside that window those records are lost, and for evidence that
is the same as their never having existed. So batch size is not a performance
knob to be turned as far as it goes---it is a second exposure window, and it
belongs next to $\Delta$ in the conformance claim. Small for irreversible
actions, generous for internal steps, on exactly the same reasoning.

The useful conclusion is that the dominant cost in our measurements is not
inherent. Anyone who reads Table~\ref{tab:latency}, concludes ``signing is too
expensive,'' and starts redesigning should read Figure~\ref{fig:batch} first.

\subsection{On real hardware, at last}
\label{sec:tdx}

Everything above was measured with the enclave as an in-process object, and we
have been careful to say so. That stand-in reproduces the protocol but leaves
the load-bearing claim untested: that a verifier can establish \emph{which code
produced this evidence} without trusting the operator. So we ran the pipeline
inside a real Intel TDX trust domain on Google Cloud, with an identical
non-confidential instance of the same shape as a control, so that the cost of
the trust domain is isolated rather than confounded with a change of CPU.

\paragraph{The binding.} A quote on its own proves that some measured
environment exists. It says nothing about which signing key belongs to it, and
an operator who could pair an honest quote with a key held outside the domain
would defeat the whole arrangement. TDX provides 64 bytes of
\texttt{REPORTDATA} that the hardware copies into what it signs, so we put the
digest of the evidence public key there. The hardware signature then covers the
measurement \emph{and} the key together, and the statement a verifier receives
is the one that matters---this key lives inside this measured
environment---rather than two facts that merely arrived in the same envelope.

It works. The instance reports \texttt{Intel TDX} and \texttt{tdx: Guest
detected}; a quote is 8{,}000 bytes and carries a real
MRTD (\texttt{c1ee9c16e3af\ldots}); the quote commits to our key, and rejects a
key we did not attest. This is requirement C7.2.4, and it is the first claim in
this paper established by hardware rather than by argument.

\paragraph{The cost, which is not where we expected it.}
Table~\ref{tab:tdx} has the numbers, and they separate cleanly into two very
different quantities.

\begin{table}[ht]
\centering\small
\begin{tabular}{@{}lrr@{}}
\toprule
& \textbf{Measured} & \textbf{Relative} \\
\midrule
Software pipeline, TDX off (control) & \SI{152.88}{\micro\second} & --- \\
Software pipeline, inside the TD & \SI{162.33}{\micro\second} & $+6.2\%$ \\
\midrule
\textbf{Trust-domain overhead} & \textbf{\SI{9.45}{\micro\second}} & \textbf{6.2\%} \\
\textbf{One hardware quote} & \textbf{\SI{39.5}{\milli\second}} & \textbf{4{,}180$\times$ the overhead} \\
\bottomrule
\end{tabular}
\caption{The pipeline on Intel TDX (GCP \texttt{c3-standard-4}) against an
identical non-confidential instance. Running inside the trust domain is nearly
free; attesting it is not.}
\label{tab:tdx}
\end{table}

Running \emph{inside} a trust domain costs 6.2\%---about ten microseconds a
step. That is the number we had no way to measure before and were careful not
to guess at, and it is small enough to be uninteresting, which is the best kind
of result for a design that depends on it.

Producing a \emph{quote} costs \SI{39.5}{\milli\second}. That is about four
thousand times the per-step overhead, and 2.6 times the entire
\SI{15}{\milli\second} budget for a single action. Attesting once per action is
therefore not expensive but impossible, and no amount of engineering will close
a gap of that shape. Every real deployment will take one quote and use it for
many actions: amortized over 100 actions the cost falls to
\SI{557}{\micro\second} a step, and over 1{,}000 to \SI{202}{\micro\second},
which is back inside the budget.

\paragraph{Which creates an exposure window we had not named.} If a quote is
taken once and covers the next thousand actions, then those actions rest on a
measurement made before them. Should the measured code change in between, their
evidence attests an environment that is no longer the one that ran, and nothing
in the record says so.

This is the same shape as the anchoring interval $\Delta$ of
\S\ref{sec:deltameasured}: a parameter that trades cost against how much can
happen unobserved, where the only wrong answer is not to state it. So it
becomes a requirement in the same form. C7.2.3 obliges a conformance claim to
declare a maximum attestation refresh interval and to refresh within it or stop.
We would not have thought to write that requirement from the software model,
because in the software model a measurement is free and staleness never arises.

\paragraph{What this still does not establish.} The reviewers of this paper
asked for a real TEE, and this is one, but it answers a narrower question than
the one they were asking. It shows the protocol runs on real hardware, what the
hardware costs, and that the key binding holds. It does \emph{not} demonstrate
resistance to a hostile host: we did not compromise a hypervisor, and Google
operates the one underneath this. The adversary of \S\ref{sec:security} owns the
OS; what we have shown is that the construction works on hardware that claims to
resist such an adversary, not that the resistance holds.

\paragraph{And who you must still trust.} The most useful output of running on
real hardware was the list it forced us to write down. To believe a measurement
from this deployment, a verifier accepts: Intel, that the TDX module and CPU
behave as specified and the provisioning key is not misissued; Intel's
certification service, that the collateral fetched to check the quote is
authentic; the quoting enclave, that it signs reports honestly; Google, that the
host does not extract domain memory by some route outside the TDX threat model,
and that the firmware measured into the runtime registers is what it says;
and whoever publishes reference values, that the MRTD being compared against
belongs to the code one believes it does.

Anchoring the evidence log removes \emph{none} of these. It makes the
\emph{history} publicly auditable, which is a real and separate property. The
\emph{measurement} is still established by that chain of parties. Our tier
definitions said that at Tier 3 there is nobody left to trust, and this list is
the plain refutation of that sentence. The standard has since been corrected
against it---the Tier~3 column no longer claims the trusted party is removed---%
and we take up what that changes, and what it does not, in
\S\ref{sec:discussion}.

\subsection{The question a chain cannot answer}
\label{sec:merkle}

Everything so far treats verification as an all-or-nothing act: you take the
history and you replay it. That is fine when the history is a few thousand
records. It stops being fine the moment you ask the question a real auditor
asks, which is not ``is the whole log intact?'' but ``was \emph{this} action in
the log?''

With a hash chain there is no way to answer the second question except by
answering the first. Each link depends on the one before, so to convince
yourself that record 500{,}000 is genuine you fetch all million records and
rehash them. An auditor spot-checking one action has to download the entire
log. That is not an inefficiency you optimize away later; it decides whether
audit is something you do continuously or something you do once a year with a
data-transfer budget and a sense of dread.

The fix is not ours and it is not new. Certificate Transparency solved this in
2013 with an append-only Merkle tree~\cite{rfc6962}, and we have been citing CT
throughout this paper while quietly using a structure that gives up its main
advantage. So we built the tree and measured what it buys.

\paragraph{Checking one record.} Table~\ref{tab:merkle} is the whole argument.
An inclusion proof at 100{,}000 records is 544 bytes and takes
\SI{7}{\micro\second} to check. Replaying the chain to learn the same fact
takes \SI{70}{\milli\second} and \SI{123}{\mega\byte} of transfer---about
69{,}000 times more data (Figure~\ref{fig:merkle}). And the gap is not fixed:
one curve is linear and the other is logarithmic, so every record you add makes
the comparison worse.

\begin{table}[ht]
\centering\small
\begin{tabular}{@{}rrrrrr@{}}
\toprule
\textbf{Records} & \textbf{Replay} & \textbf{Check a proof} &
\textbf{Speedup} & \textbf{Proof} & \textbf{Fetched to replay} \\
\midrule
10 & \SI{8}{\micro\second} & \SI{1.87}{\micro\second} & 4$\times$ & 128 B & \SI{12}{\kilo\byte} \\
1{,}000 & \SI{672}{\micro\second} & \SI{4.20}{\micro\second} & 160$\times$ & 320 B & \SI{1.2}{\mega\byte} \\
10{,}000 & \SI{6.9}{\milli\second} & \SI{5.67}{\micro\second} & 1{,}217$\times$ & 448 B & \SI{12.3}{\mega\byte} \\
100{,}000 & \SI{70}{\milli\second} & \SI{7.01}{\micro\second} & \textbf{9{,}972$\times$} & 544 B & \SI{122.7}{\mega\byte} \\
\bottomrule
\end{tabular}
\caption{Establishing that one named record is in the log. The chain must be
replayed from the beginning; the tree needs a proof of $\lceil \log_2 n \rceil$
hashes.}
\label{tab:merkle}
\end{table}

\begin{figure}[t]
\centering
\includegraphics[width=0.88\linewidth]{figures/chart-merkle.pdf}
\caption{What an auditor must download to check a single action. One line is
linear in log size and the other is logarithmic, which is why the gap widens
with every record written.}
\label{fig:merkle}
\end{figure}

Put it in operational terms. A thousand agents at a hundred actions a day
produce 36.5 million records a year, a tree 26 levels deep. Checking one
sampled action costs about 2 kilobytes. Replaying the chain costs
\SI{44.8}{\giga\byte}---and costs the same whether you wanted to check one
record or a million, because replay is all-or-nothing. Proofs cost what you
actually asked for.

\paragraph{What it costs the agent.} A structure that helps the auditor at the
agent's expense would be a bad bargain, so this is the number that decides
adoption. Maintaining the tree adds \SI{2.1}{\micro\second} per action at ten
records and \SI{3.2}{\micro\second} at a hundred thousand---between 2.4\% and
3.8\% of the \SI{85}{\micro\second} a signature already costs. It is close to
free, and it stays close to free as the log grows, because appending touches
$\log_2 n$ nodes rather than $n$.

\paragraph{The part where it broke our own system.} We expected this experiment
to produce a performance table. It produced attack A9 instead, and A9 is the
most useful thing in this section.

Every attack in Table~\ref{tab:attacks} except this one assumes the adversary
cannot sign. But the operator running the enclave \emph{can}---that is what
operator-controlled infrastructure means, and it is the threat model we claim
to be working in. Give an operator the key and it can reach into the middle of
the log, change a record, recompute every link after it, and re-sign the whole
thing. What comes out has exactly the right length. Every signature verifies.
Every link matches. Replaying it finds nothing wrong, because there \emph{is}
nothing internally wrong with it: it is a perfectly consistent history of
events that did not happen.

The only thing in the world that contradicts it is a root published before the
rewrite. Our verifier had one---and compared the published \emph{step count}
while never comparing the published \emph{root}. So it accepted the forgery,
cheerfully, and told us the chain verified.

We had written the truncation check and stopped thinking. Truncation is the
attack where the history is too short, so we checked the length, and checking
the length felt like checking the anchor. It is not. Comparing roots catches
the rewrite immediately, and a consistency proof does better still: it settles
in $O(\log n)$ not merely \emph{that} two published roots disagree but that the
newer history is not an extension of the older one---which is the precise
statement that history was rewritten. Those proofs run 224 to 480 bytes and
verify in under \SI{9}{\micro\second}, and in our harness they catch every
rewrite we could construct.

This is now requirement C7.3.5, and its wording---``compares the recomputed
root against the anchored value rather than only against a step count''---is
oddly specific for a standard because it is a description of our own mistake.
We would rather the requirement be strange and correct than elegant and
silent about the thing that actually goes wrong.

There is a general lesson here that we did not expect to be writing down. Two
earlier findings in this paper came from proving theorems; this one came from
building a component we thought was an optimization. Implementing the fast path
forced us to state precisely what a published root was supposed to guarantee,
and stating it precisely was enough to notice that we were not checking it. The
performance work found a security bug, which is not the direction that traffic
usually runs.

\subsection{Remaining parameters}
For the \textbf{privacy path}, an implementation admitting zero-knowledge
evidence must state the relation proven, the setup assumption, and what the
disclosed statement leaks. ``We use ZKPs'' is not a claim anyone can check---%
and if the setup was run by one party, you have quietly reintroduced exactly
the trusted party the whole scheme exists to remove.

\section{What We Still Don't Know}
\label{sec:discussion}

Standards-integration targets and the phased rollout are in
Appendices~\ref{sec:sdo} and~\ref{sec:roadmapapp}. This section is for the
things we have not settled.

\subsection{The honest list}
\textbf{(1)~The enclave is no longer pretend, and the honest limit has moved.}
An earlier draft of this list said the hardware was not there and that
measuring on TDX or SEV-SNP was the obvious next job. It has since been done:
\S\ref{sec:tdx} reports the pipeline running inside a real Intel TDX trust
domain on Google Cloud, against an identical non-confidential instance as a
control. Running inside the trust domain costs \SI{9.45}{\micro\second} a
step, 6.2\%---the number we had refused to guess at, and small enough to be
uninteresting. Producing a quote costs \SI{39.5}{\milli\second}, which is 2.6
times the entire per-action budget and made per-action attestation not
expensive but impossible; that is where requirement C7.2.3 came from, and we
would not have thought to write it from the software model, where a measurement
is free and staleness never arises.

What remains unmeasured is narrower and should not be read as settled. We did
not compromise a hypervisor, and Google operates the one underneath this
deployment, so what we show is that the construction works on hardware that
claims to resist a hostile host---not that the resistance holds. The adversary
of \S\ref{sec:security} owns the OS, and that adversary has not been run
against this. SEV-SNP and ARM CCA are unmeasured entirely, and a second
platform is the obvious next job now.

\paragraph{Our tier definitions overclaimed, and we now have the receipts.} The
standard said that at Tier 3 there is nobody left to trust. \S\ref{sec:tdx}
lists five parties a verifier must still accept to believe a measurement from
our own TDX deployment, and anchoring removes none of them. The sentence
conflates two properties that a careful reader will separate for us: evidence
can be \emph{publicly verifiable}---anyone may check it with published
tools---without being \emph{trust-independent}, and hardware attestation buys
the first while leaving a vendor chain intact underneath the second.

That sentence is now gone. The Tier~3 column says the evidence is produced by
the mechanism and checkable by anyone with published tools, and that the parties
which remain load-bearing are disclosed rather than removed; the word
``trustless'' has been struck from the normative text, where it occurred once.
What has \emph{not} been settled is the framing that replaces it---whether the
Tier is demoted from the primary claim to a coarse summary, with what must be
trusted carried by the trust-assumption disclosure instead, and whether the
Tier-2/Tier-3 test is restated as who selects the trusted party and whether
their dishonesty is publicly detectable. Those are working-group decisions and
are marked as such.

We note what the correction rests on, because it bounds it: five self-authored
deployment descriptions and no real-world system. That is enough to refute a
universal claim---one counterexample does it, and there are two disjoint
ones---and it is not enough to install a new ordering or a new taxonomy in its
place. We would rather print that here than let it be discovered.

\textbf{(2)~We verify what happened, not whether it was wise.} No claim about
correctness, fairness, or intent. An agent that does something harmful
entirely within its permissions is a safety problem, and we do not solve
safety problems.

\textbf{(3)~Trust is reduced, not abolished.} Everything above Tier~2 rests on
(A1), (B1), (C1). Enclaves have been broken~\cite{costan2016,foreshadow}, and
the chip vendor is a party you are trusting. Our answer is to make you write
that down, not to pretend it is not there.

\textbf{(4)~The framework comparison is one coder and one direction}
(Section~\ref{sec:gaplimits}). Treat it as a sketch.

\textbf{(5)~The hard line may move.} Where exactly Tier~2 ends and Tier~3
begins is under review by people who do cryptography and complexity theory for
a living, including work on what is provably impossible to verify efficiently.
If they tell us the line is in the wrong place, the line moves.

\textbf{(6)~Failing closed means failing.} At Tier~4, if evidence cannot be
produced, work stops. That is the property we want and also an availability
risk, and both facts have to be disclosed. You do not get to advertise only
the first one.

Open questions the working group is still arguing about---who owns identity
binding, whether unlinkable-but-verifiable identity should be allowed, whether
evidence continuity across jurisdictions is a fifth property, and how
continuous ``continuous'' has to be---are tracked publicly.

\subsection{Where this stands}
\poc{} is a working draft under public comment, with a public repository
(\url{https://github.com/AAI-Society/ov-poc-standard}) containing the
specification, the audit checklist, the coding sheets, and the
implementation with its attack harness and benchmarks. We would rather have
this attacked than admired. The most useful thing anyone could do is run the
utility measurement on real agent traces instead of our synthetic ones,
because that number---not the cryptography---is what will decide whether this
is deployable.

\section{Conclusion}

The way we have always known what software did is that we asked it and it told
us. That worked when software followed instructions, and it does not work now
that software makes decisions, at machine speed, in situations nobody
enumerated in advance. Asking an agent what it did is not a check on the
agent; it is a conversation with the party you would be checking.

What we have described is the alternative: evidence made by the mechanism
while it acts, chained so gaps show, published so nobody can keep two sets of
books, tied to the effect so it cannot describe a fiction, and gradable so a
buyer can ask one yes-or-no question. We proved the parts that are provable,
and we were careful to say which parts are not: an attester by itself proves
less than it appears to, and a monitor that watches the path will refuse half
your work unless you tell it where to forget.

We also built it, which is how we learned that the cryptography costs
\SI{201}{\micro\second} and is not the problem, that the policy design costs
42\% of your legitimate work and is, and that making the whole thing survive
quantum computers is a storage bill rather than a redesign. Then we ran it on
real Intel TDX hardware, which settled the question the software model could
not: running inside a trust domain costs 6.2\%, about ten microseconds a step,
and is not where the expense lives. Producing one hardware quote costs
\SI{39.5}{\milli\second}---2.6 times the entire budget for a single
action---so attesting per action is not expensive but impossible, and every
deployment will amortize one quote across many. That is an exposure window
rather than an optimization, and naming it became a requirement we had no
reason to write while the enclave was a stand-in. Building it is also
how we found the bug in our own verifier, which no amount of rereading the
specification had turned up: we compared a published step count where we should
have compared a published root, and a rewritten history sailed straight through.
We would not have noticed it by thinking harder. We noticed it by writing the
code that made us say precisely what a published root was for.

The specification, the code, the attacks, the schema, the test vectors, and the
numbers are all open. If any of it is wrong, it should be possible for you to
find out---which, in the end, is the entire point.


\section*{Acknowledgments}
This work is developed within the \poc{} Initiative of the Advanced AI
Society. We thank the domain working groups and contributors whose input
shaped the draft standard, including Bob Blessing-Hartley, Advait Patel,
David Thomson, Drummond Reed, and Hart Montgomery (Linux Foundation
Decentralized Trust), and the Cloud Security Alliance MAESTRO and AARM
communities. The \poc{} Lab is being established as a community lab at Linux
Foundation Decentralized Trust.

\appendix

\section{Standards Integration}
\label{sec:sdo}
\emph{ISO/IEC JTC 1/SC 42:} \poc{} provides a technical realization for
active SC~42 workstreams~\cite{jtc1sc42,nemko2025,lexai2025}---a technical
annex for auditable runtime controls under ISO/IEC~42001 Clauses~8--9
(supplying cryptographically verifiable compliance evidence during
certification audits)~\cite{iso42001}, operationalization of ISO/IEC~23894
risk treatment as executable, hardware-verified policy
bundles~\cite{iso23894}, and an SC~42/SC~27 bridge coupling AI oversight
with confidential-computing standards~\cite{iec2025}.
\emph{IETF RATS:} an Internet-Draft establishing an AI-agent EAT
profile~\cite{rfc9711,tcg2024}---standardizing claims for lifecycle hook
types, policy-bundle hashes, Agent Bill of Materials (AgBOM) digests, and
step indices---and an extension of multi-verifier attestation to sub-agent
delegation chains~\cite{multiverifier}.
\emph{ACS:} \poc{} supplies the missing cryptographic layer for
ACS~\cite{acs2026,acsms2026}, anchoring its decision engine in TEEs and
signing decision tokens.
\emph{NIST AI 600-1:} \poc{} provides concrete implementations for the
generative-AI profile's measurement
controls~\cite{nist6001,ciphernorth2026,modulos2026}---enforcing that
deployment configurations cannot bypass security controls and satisfying
continuous security verification with attestation tokens for every tool
call and state transition.

\section{Implementation Roadmap}
\label{sec:roadmapapp}
Figure~\ref{fig:roadmap} summarizes the path to ratification and adoption.
\textbf{Phase 1 (months 1--6): open specification and reference core}---core
schema extensions on the ACS repository, open-source Rust reference enclave
policy engines, and the IETF Internet-Draft for the AI-agent EAT profile.
\textbf{Phase 2 (months 7--12): multi-vendor interoperability and SDO
balloting}---SC~42 working groups, conformance testing across Python,
Node.js, and Rust agent frameworks, and telemetry integration with
OpenTelemetry and OCSF. \textbf{Phase 3 (months 13--24): ratification and
deployment}---normative annexes in ISO/IEC~42001 and 23894, RFC publication,
and deployments across financial, healthcare, and public-sector workloads.
On the adopter side, the same three phases carry an organization from
foundational key, signing, and logging infrastructure through expanded
proof coverage to continuously monitored operation---each phase reaching
readiness for one conformance stage.

\begin{figure}[t]
\centering
\includegraphics[width=0.98\linewidth]{figures/roadmap.pdf}
\caption{Implementation roadmap. Each phase brings an adopting organization
to readiness for one conformance stage; on the standards track, Version 1.0
targets February 1, 2027.}
\label{fig:roadmap}
\end{figure}


\bibliographystyle{unsrt}
\bibliography{references}

\end{document}
